\chapter{系统设计与方法实现}

围绕"基于大模型的软件运维方法"这一研究主题，本文设计并实现了一个面向软件运维场景的智能故障诊断原型系统——LLM-Ops Copilot。本章从系统设计目标出发，依次介绍领域知识库构建、证据融合方法、RCA推理引擎和原型系统实现。

\section{系统设计目标}

本文原型系统的设计目标包括以下四个方面。

第一，知识增强能力。构建面向运维领域的知识库，通过检索增强生成机制为模型提供标准操作流程、历史故障案例等背景信息，提升回答的专业性。

第二，多源数据组织能力。将日志、监控指标和告警等异构数据统一组织为适合大模型分析的上下文，支撑后续推理任务。

第三，故障诊断能力。通过结构化提示引导大语言模型完成故障摘要生成、原因推断与修复建议输出。

第四，自然语言交互能力。基于ChatOps理念提供对话式交互界面，降低运维人员的使用门槛。

\section{运维知识库构建}

本文构建的运维知识库涵盖标准操作流程（SOP）文档、历史故障案例和系统架构说明等。其中SOP文档包含CPU故障、内存故障、数据库故障、网络故障等常见场景的标准排查步骤和修复方案。

知识库构建过程包括四个步骤。第一步，对原始文档进行格式清理，去除无关符号和重复内容。第二步，按段落和标题层级对长文档进行切分，每个片段保留必要的上下文信息。第三步，使用sentence-transformers预训练模型将文本片段转换为768维语义向量~\cite{zhang2024rag4itops}。第四步，将向量与原始文本、元数据（服务名、故障类型等）一同存入ChromaDB向量数据库，支持按服务名和故障类型进行元数据过滤检索~\cite{zhang2024rag4itops}。

知识检索时，系统根据用户问题生成查询向量，在向量数据库中检索语义最相似的Top-K知识片段，并按相关性排序后返回给后续模块。

\section{多源数据融合方法}

在故障诊断场景中，单一数据源往往难以支撑高质量分析~\cite{pham2024rcacausal}。本文设计了证据融合器（EvidenceBuilder）模块，将知识检索结果、日志摘要和监控指标描述统一组织为结构化的分析上下文。

证据融合器的处理流程如下。首先，对日志数据进行筛选和摘要：按服务名和时间范围过滤日志，提取错误日志和警告日志，统计各类关键词出现频次，将原始日志信息浓缩为结构化的日志摘要。其次，对监控指标进行处理：按预设阈值识别异常指标，将数值型指标异常信息转换为文本描述。最后，将知识检索结果、日志摘要和指标摘要组织为统一的AnalysisContext数据结构。

在严重程度评估方面，证据融合器综合考虑以下因素：异常指标数量越多、错误日志越多、指标峰值越高，则故障越严重。该评估结果作为元信息传递给后续的RCA引擎，帮助模型判断故障的紧迫程度。

证据融合的完整流程如算法~\ref{algo:evidence-builder}所示。

\begin{algorithm}
\caption{证据融合算法}
\label{algo:evidence-builder}
\DontPrintSemicolon
\SetKwInput{KwIn}{输入}
\SetKwInput{KwOut}{输出}
\KwIn{$question$, $service$, $time\_range$, $knowledge\_results$, $metrics$, $logs$}
\KwOut{$context$}
\nl $anomalies$ $\gets$ 从$metrics$中提取异常指标\;
\nl $severity$ $\gets$ 根据异常数量和日志情况评估严重程度\;
\nl $context$ $\gets$ 构建AnalysisContext结构，汇总知识检索、日志摘要、指标摘要和严重程度\;
\nl \Return{$context$}
\end{algorithm}

\section{基于大模型的根因分析引擎}

根因分析（RCA）是本文方法的重点任务~\cite{roy2024agents}。本文设计并实现了RCA推理引擎（RCAEngine），通过结构化Prompt引导大语言模型完成故障分析。

本文设计的RCA Prompt包含三个核心部分。

第一，系统级角色设定。定义模型为资深SRE（网站可靠性工程师），明确其分析原则为：证据驱动（所有结论必须基于提供的证据）、结构化输出（严格按格式返回）、置信度评估（每个根因需给出置信度）和诚实透明（证据不足时明确说明）~\cite{liu2023prompt,liu2023logprompt}。

第二，上下文输入规范。将知识检索结果、日志摘要、指标摘要和故障严重程度等信息结构化地组织为Prompt上下文~\cite{zhang2024rag4itops}。

第三，推理步骤约束。引导模型按照以下步骤展开分析：（1）现象总结——简洁描述当前异常表现；（2）证据分析——从指标、日志和知识库中提取关键信息；（3）根因假设——列出1至3个可能的根因，并给出置信度和支撑/反对证据；（4）验证计划——列出验证每个根因的具体检查步骤；（5）修复建议——针对最可能的根因给出可操作的修复方案~\cite{roy2024agents}。

RCA引擎还设计了后备规则分析机制。当大模型调用失败（如超时或解析错误）时，系统自动切换到基于规则的简单分析~\cite{zhang2024rag4itops}：根据异常指标类型推断根因（如CPU异常推断为CPU密集型任务，连接异常推断为连接池问题），确保系统在任何情况下都能返回一定形式的分析结果。

RCA推理的完整流程如算法~\ref{algo:rca-engine}所示。

\begin{algorithm}
\caption{RCA推理算法}
\label{algo:rca-engine}
\DontPrintSemicolon
\SetKwInput{KwIn}{输入}
\SetKwInput{KwOut}{输出}
\KwIn{$context$}
\KwOut{$result$}
\nl $prompt$ $\gets$ 根据上下文构建结构化Prompt\;
\nl $raw\_response$ $\gets$ 调用大语言模型API\;
\nl \If{解析成功}{
    \nl \Return{结构化RCAResult}
}
\nl \Else{
    \nl $result$ $\gets$ 调用后备规则分析\;
    \nl \Return{$result$}
}
\end{algorithm}

\section{原型系统实现}

原型系统采用前后端分离架构。前端基于Streamlit框架构建，提供故障诊断、知识查询、监控展示和实验评估四个人机交互入口。用户通过自然语言向系统提问，系统返回包含根因分析、置信度和修复建议的结构化诊断结果。

后端基于FastAPI框架实现，提供分析、知识检索和实验评估等RESTful接口。后端整合了知识检索服务、日志服务、指标服务和模型调用服务等模块~\cite{zhang2024rag4itops}。

端到端的RAG Pipeline执行流程为：接收用户问题后，首先调用知识检索服务获取相关SOP文档和历史案例；然后调用日志和指标服务获取实时故障信息；接着通过证据融合器将多源信息统一组织为分析上下文；最后调用RCA推理引擎生成诊断结果并返回给用户。

\section{本章小结}

本章介绍了LLM-Ops Copilot原型系统的设计与实现。系统通过构建面向运维领域的知识库，利用检索增强生成技术获取领域知识背景；通过证据融合器将日志摘要、监控指标和知识检索结果统一组织为结构化上下文；通过结构化Prompt引导大语言模型按照"证据分析—根因假设—验证计划—修复建议"的逻辑顺序展开诊断分析。原型系统在技术层面验证了上述方法在运维故障诊断任务中的可行性，也为后续实验评估奠定了实现基础。
